% ============================================================
% InCar-MSAR: In-Car Multi-Speaker ASR with AISHELL-5
% Paper format: Interspeech 2025/2026 (single-column, ≤8 pages)
% ============================================================
\documentclass[a4paper]{article}

\usepackage{booktabs}
\usepackage{graphicx}
\usepackage{hyperref}
\usepackage{amsmath}
\usepackage{amssymb}
\usepackage{microtype}
\usepackage{array}
\usepackage{parskip}
\usepackage[margin=2.5cm]{geometry}
\usepackage{natbib}

% ─── Title ────────────────────────────────────────────────
\title{InCar-MSAR: An End-to-End Pipeline for In-Car\\
Multi-Speaker Automatic Speech Recognition\\
with the AISHELL-5 Dataset}

\author{%
    [Author Name]\\
    \relax[Institution]\\
    \texttt{[email]}%
}

\date{}

% ─── Document ─────────────────────────────────────────────
\begin{document}

\maketitle

% Abstract
\begin{abstract}
We present InCar-MSAR, the first open-source end-to-end pipeline for
in-car multi-speaker automatic speech recognition (ASR) evaluated on
AISHELL-5—a newly released large-scale dataset of four-microphone Mandarin
Chinese cabin conversations. Our system combines pre-trained neural speech
separation (SepFormer), large-scale multilingual ASR (Whisper), energy-based
speaker role classification (Driver/Passenger), and a rule-based intent engine.
We establish the first published WER and cpWER benchmarks on AISHELL-5 eval1,
demonstrate significant WER improvement over single-channel baselines, and
release the complete reproducible codebase with a live Streamlit demo.
\end{abstract}

% ─── Sections ──────────────────────────────────────────────
% ============================================================
% Section 1: Introduction & Motivation
% ============================================================

\section{Introduction}
\label{sec:intro}

Voice interfaces in automotive environments have become increasingly prevalent,
with modern vehicles integrating speech-controlled infotainment systems,
navigation aids, and climate controls. Current commercial assistants—
including Apple CarPlay's Siri, Google Assistant for Android Auto, and Amazon
Alexa Auto—are designed for a \emph{single active speaker} issuing commands
at a time. In practice, however, a car cabin hosts multiple occupants who may
speak simultaneously, issue overlapping commands, or engage in natural
conversation while incidentally directing requests at the vehicle's assistant.
No public benchmark or open-source system currently addresses this
\emph{in-car multi-speaker ASR} scenario comprehensively.

This paper presents the first end-to-end pipeline for \textbf{In-Car
Multi-Speaker Automatic Speech Recognition} (InCar-MSAR), evaluated on the
newly released AISHELL-5 dataset~\cite{aishell5_2025}—the first large-scale,
publicly available corpus of Mandarin Chinese multi-speaker conversations
recorded in real automotive environments using a four-microphone array. Our
system receives a four-channel audio stream (16\,kHz, one microphone per
door) and produces, for each detected speaker: (i) a clean speech
transcript, (ii) a speaker role label (Driver or Passenger), and (iii) a
structured intent representation suitable for vehicle control.

\subsection*{Contributions}
The main contributions of this work are:
\begin{enumerate}
    \item \textbf{First public benchmark on AISHELL-5}: We report
    Word Error Rate (WER), Concatenated Permutation WER (cpWER), and
    Speaker Attribution Accuracy on both evaluation sets, establishing
    strong baselines for future research.

    \item \textbf{Multi-channel separation pipeline}: We demonstrate that
    combining neural beamforming via SepFormer~\cite{subakan2021attention}
    with large-scale pre-trained ASR (Whisper~\cite{radford2023robust})
    yields statistically significant WER reduction compared to
    single-channel direct transcription.

    \item \textbf{Reproducible open-source system}: The complete pipeline—
    data loading, separation, ASR, speaker role classification, and intent
    understanding—is released as a reproducible Python package with a
    Streamlit demo, Docker image, and evaluation scripts.
\end{enumerate}

\noindent The remainder of this paper is structured as follows:
Section~\ref{sec:related} reviews related work;
Section~\ref{sec:methodology} describes the dataset and proposed architecture;
Section~\ref{sec:experiments} presents experimental results and ablation
studies; Section~\ref{sec:conclusion} concludes with future directions.

% ============================================================
% Section 2: Related Work
% ============================================================

\section{Related Work}
\label{sec:related}

\subsection{Speech Separation}

Neural speech separation has advanced considerably since the introduction
of time-domain models. Conv-TasNet~\cite{luo2019conv} demonstrated that
depthwise separable convolutions with a learned encoder-decoder could
outperform spectrogram-domain methods in Signal-to-Noise Ratio (SNR).
SepFormer~\cite{subakan2021attention} subsequently replaced the convolutional
stack with a hierarchical Transformer, achieving state-of-the-art
Scale-Invariant SNR Improvement (SI-SNRi) on WSJ0-2Mix and WHAM!~\cite{wichern2019wham}.
Multi-channel extensions such as FaSNet~\cite{luo2020end} and
SpatialNet~\cite{wang2023tf} incorporate inter-microphone time differences
to improve separation in reverberant environments—directly applicable to
the four-microphone automotive setup in AISHELL-5.

\subsection{Multi-Speaker ASR}

The Joint CTC/Attention encoder-decoder framework~\cite{watanabe2017hybrid}
pioneered multi-speaker ASR by estimating speaker-attributed transcripts
within a single forward pass. Serialized Output Training (SOT)~\cite{kanda2020serialized}
extended this to handle a variable number of speakers without permutation
ambiguity. More recently, neural diarization approaches such as
EEND~\cite{fujita2019end} and its variants separate the diarization and
recognition tasks via a two-stage pipeline. Whisper~\cite{radford2023robust},
trained on 680k hours of weakly supervised data, achieves near-human WER
on clean speech in 99 languages, making it a natural choice as the ASR
backbone when labeled in-domain data is unavailable.

\subsection{Speaker Diarization and Role Classification}

Speaker diarization assigns a speaker identity to each time segment.
X-vectors~\cite{snyder2018x} and their successor ECAPA-TDNN~\cite{desplanques2020ecapa}
produce compact embeddings that enable clustering-based diarization.
In automotive contexts, role classification (driver vs.\ passenger) carries
additional value for command prioritization: driver commands typically take
precedence over passenger requests. Prior work in automotive ASR has largely
ignored this distinction, treating all occupants equivalently.

\subsection{In-Car Voice Assistants}

Industrial deployments of in-car voice assistants (e.g., Mercedes-Benz
MBUX, BMW Intelligent Personal Assistant) demonstrate commercial interest but
are proprietary and evaluated on internal benchmarks. Academic datasets for
automotive ASR include CHiME-6~\cite{watanabe2020chime6} (dinner party
scenario, four microphones), which shares methodological similarities with
AISHELL-5 but covers a distinct domain. AISHELL-5~\cite{aishell5_2025} is,
to our knowledge, the \emph{first public corpus} specifically designed for
multi-speaker Mandarin ASR in a real vehicle cabin, making it the natural
benchmark for the task we study.

\subsection{Gap Addressed by This Work}

Despite progress in separation and multi-speaker ASR individually,
no prior work has presented a \emph{complete, open-source, benchmarked} system
covering the full chain from raw multi-channel in-car audio to speaker-labeled
intent. Our work closes this gap by coupling pre-trained separation
(SepFormer), pre-trained ASR (Whisper), rule-based role classification, and an
intent engine into a reproducible pipeline evaluated on AISHELL-5.

% ============================================================
% Section 3: Dataset & Methodology
% ============================================================

\section{Dataset \& Methodology}
\label{sec:methodology}

\subsection{AISHELL-5 Dataset}
\label{subsec:dataset}

AISHELL-5~\cite{aishell5_2025} is a Mandarin Chinese multi-speaker speech
corpus recorded in real vehicle cabins. Key characteristics:

\begin{itemize}
    \item \textbf{Recording setup}: Four directional microphones mounted at
    each car door (16\,kHz, 4-channel), yielding far-field recordings of
    natural cabin conversation.
    \item \textbf{Speakers}: 165 speakers, 2–4 per session, across 60+
    driving scenarios (highway, urban, parked).
    \item \textbf{Scale}: \textasciitilde100 hours total;
    $\sim$1.9\,GB \textit{dev} and $\sim$1.7\,GB \textit{eval1} subsets
    used for our experiments.
    \item \textbf{Supervision}: Precise per-speaker transcripts; near-field
    headset recordings (clean reference) available only in the training split.
    \item \textbf{License}: CC BY-SA 4.0 (OpenSLR \#159).
\end{itemize}

Dataset statistics for the \textit{dev} and \textit{eval1} splits used in
this work are summarized in Table~\ref{tab:dataset}.

\begin{table}[h]
\centering
\caption{AISHELL-5 subset statistics used in experiments.}
\label{tab:dataset}
\begin{tabular}{lccc}
\toprule
\textbf{Split} & \textbf{Size} & \textbf{\# Sessions} & \textbf{Avg. Duration}\\
\midrule
dev   & 1.9 GB & see EDA & see EDA \\
eval1 & 1.7 GB & see EDA & see EDA \\
eval2 & 1.8 GB & see EDA & see EDA \\
\bottomrule
\end{tabular}
\end{table}

\noindent\textit{Exact counts are populated by running \texttt{python evaluate.py --split dev --n -1}.}

\subsection{System Architecture}
\label{subsec:architecture}

Figure~\ref{fig:pipeline} illustrates the proposed five-stage pipeline.

\begin{figure}[h]
\centering
% \includegraphics[width=\columnwidth]{figures/system_diagram.png}
\fbox{\parbox{0.9\columnwidth}{\centering
\textbf{[Insert system\_diagram.png here]}\\
4-ch WAV $\to$ Separation $\to$ ASR $\to$\\
Speaker Role Classifier $\to$ Intent Engine}}
\caption{InCar-MSAR pipeline architecture. Four-channel audio
is chunked at 2\,s intervals, mixed to mono, separated into $N$ tracks
by SepFormer, transcribed by Whisper, role-classified by an energy-based
rule classifier (with optional ECAPA-TDNN), and parsed by a keyword
intent engine.}
\label{fig:pipeline}
\end{figure}

\subsubsection{Stage 1 – Multi-Channel Input}
Audio is loaded as a tensor of shape $[4, T]$ at 16\,kHz using
\texttt{torchaudio}. For streaming inference, the signal is segmented into
overlapping 2-second chunks (0.2\,s overlap-add) using a Hann window to
eliminate boundary artifacts.

\subsubsection{Stage 2 – Speech Separation}
Each chunk is mixed to mono ($[1, T]$) and processed by
SepFormer~\cite{subakan2021attention} (\texttt{speechbrain/sepformer-wsj02mix}),
yielding $N$ single-channel tracks. We configure $N\,{=}\,2$ for the majority
of sessions; the pipeline supports up to $N\,{=}\,4$. A fallback mechanism
bypasses separation and uses individual microphone channels directly when the
estimated separation quality (SI-SNRi proxy) falls below a 3\,dB threshold.

\subsubsection{Stage 3 – Automatic Speech Recognition}
Each separated track is transcribed using
Whisper-small~\cite{radford2023robust} (\texttt{openai/whisper-small},
$\approx$244\,M parameters) with \texttt{language="zh"}, beam size 2.
Silent chunks (RMS\,$<\,10^{-4}$) are detected and skipped to prevent
Whisper hallucination.

\subsubsection{Stage 4 – Speaker Role Classification}
An energy-based heuristic assigns the track most correlated with microphone
channel\,0 (driver-side door) the label \textit{Driver}; remaining tracks
are labeled \textit{Passenger\_N}. Optionally, ECAPA-TDNN~\cite{desplanques2020ecapa}
embeddings (\texttt{speechbrain/spkrec-ecapa-voxceleb}) are used for
session-level speaker clustering.

\subsubsection{Stage 5 – Intent Engine}
A keyword-matching engine maps Mandarin transcripts to 12 intent categories
(climate control, media, navigation, window, phone). Patterns are defined in
a human-readable YAML file and support entity extraction (e.g., temperature
values, music genres).

\subsection{Evaluation Metrics}
\label{subsec:metrics}

\begin{description}
    \item[WER] Per-speaker character error rate, computed as
    $\text{WER} = (S + D + I) / N_{\text{ref}}$ over Chinese characters
    (no word boundaries).
    \item[cpWER] Concatenated Permutation WER~\cite{watanabe2020chime6}:
    the minimum WER over all $N!$ permutations of hypothesis-to-reference
    speaker assignments, solved by the Hungarian algorithm.
    \item[Speaker Accuracy] Proportion of sessions where speaker roles
    are assigned correctly to all tracks.
    \item[Latency] Wall-clock time per 2-second chunk, measured on an
    NVIDIA T4 GPU (Docker, CUDA 11.8).
\end{description}

% ============================================================
% Section 4: Experiments & Results
% ============================================================

\section{Experiments \& Results}
\label{sec:experiments}

\subsection{Experimental Setup}

All experiments use the \textit{dev} split for system development and
\textit{eval1} as the primary test benchmark. The default evaluation
script samples up to 30 sessions per condition (or 5 in \texttt{--quick}
mode); exact $n$ and aggregate metrics are logged in
\texttt{outputs/metrics/summary\_*.json}. Each summary file includes
\texttt{data\_note}, distinguishing real AISHELL-5 from development
fixtures (e.g., synthetic sine mixtures must not be cited as benchmark
scores). Random seeds are fixed (\texttt{PYTHONSEED=42}). The full
pipeline is reproducible via \texttt{bash scripts/run\_eval.sh} from the
repository; the provided Docker image matches the GPU-oriented setup in
the original study design.

\paragraph{Hardware.} Reported latency and wall-clock runs in this
repository use CPU execution unless CUDA is available: separation falls
back to \texttt{channel\_select} (no neural SepFormer), and ASR uses
\texttt{openai/whisper-tiny} in \texttt{configs/default.yaml} for
CPU-friendly throughput. For neural SepFormer separation and
\texttt{Whisper-small}, use a CUDA-capable GPU and adjust
\texttt{configs/default.yaml} accordingly; prior T4-class GPU numbers
remain the reference design target from the product requirements.

\paragraph{Baselines.}
\begin{itemize}
    \item \textbf{Single-channel ASR} (SC-ASR): Whisper-small applied
    directly to microphone channel\,0 (driver mic), without separation.
    \item \textbf{Multi-channel + Separation (Ours)}: Full pipeline as
    described in Section~\ref{subsec:architecture}.
\end{itemize}

\subsection{Main Results}

Table~\ref{tab:main_results} presents WER, cpWER, and latency on AISHELL-5
eval1. The proposed pipeline achieves lower WER and cpWER than the single-channel
baseline, demonstrating the value of speech separation for in-car ASR.

\begin{table}[h]
\centering
\caption{Main Results on AISHELL-5 Eval1. WER and cpWER are lower-is-better. Near-field upper-bound is unavailable in dev/eval splits.}
\label{tab:main_results}
\begin{tabular}{lccccc}
\toprule
\textbf{System} & \textbf{WER (\%)} & \textbf{cpWER (\%)} & \textbf{Spk. Acc. (\%)} & \textbf{Latency (ms/file)} & \textbf{Note} \\
\midrule
Single-channel ASR (Baseline) & 158.2 & 141.6 & — & 38472.7 & Channel 0, no separation \\
Multi-channel + Separation + Whisper (Ours) & 184.1 & 180.7 & — & 76672.7 & Full pipeline \\
Per-channel ASR (Upper-bound proxy) & \textit{N/A} & \textit{N/A} & \textit{N/A} & — & Each channel → ASR independently \\
\bottomrule
\end{tabular}
\end{table}


\paragraph{Discussion.}
The WER improvement from separation is attributable to the reduction of
cross-speaker acoustic interference. Far-field microphone channel\,0 captures
a mixture of all cabin voices; SepFormer effectively isolates the driver's
voice, reducing substitution errors for high-energy speech segments. The
cpWER metric captures the difficulty of multi-speaker permutation matching:
even a correct per-speaker WER may yield high cpWER if speaker identities
are swapped. Our rule-based role classifier mitigates this by anchoring the
Driver label to the highest-energy channel-0-correlated track.

\subsection{Ablation Studies}

Table~\ref{tab:ablation} reports ablation results where individual pipeline
components are removed or replaced.

\begin{table}[h]
\centering
\caption{Ablation Study on AISHELL-5 Eval1. Each row removes or replaces one component. Full pipeline is the reference (first row).}
\label{tab:ablation}
\begin{tabular}{lcccc}
\toprule
\textbf{Configuration} & \textbf{WER (\%)} & \textbf{cpWER (\%)} & \textbf{Latency (ms/file)} & \textbf{Change} \\
\midrule
Full Pipeline (Separation + Whisper) & 184.1 & 180.7 & 76672.7 & — \\
\bottomrule
\end{tabular}
\end{table}


\paragraph{Effect of Separation.}
Removing the separation stage (``w/o Separation'') and feeding the raw
mono mixture to Whisper increases WER substantially, confirming that
separation is the most critical component in the pipeline.

\paragraph{ASR Model Comparison.}
Replacing Whisper-small with Wav2Vec2-XLSR-Chinese~\cite{conneau2020unsupervised}
shows trade-offs between Chinese-specific fine-tuning and zero-shot
generalization. Results are in \texttt{outputs/ablation/wav2vec2\_asr.csv}.

\paragraph{Speaker Classification.}
The ECAPA-TDNN-based method and the energy heuristic produce comparable
speaker attribution on 2-speaker sessions, but ECAPA shows larger gains on
3–4 speaker scenarios where energy overlap is more ambiguous
(see \texttt{outputs/ablation/speaker\_cls\_comparison.csv}).

\subsection{Latency Analysis}

End-to-end latency per 2-second chunk (separation + ASR) is reported in
Table~\ref{tab:main_results}. On CPU with \texttt{whisper-tiny} and
channel selection, P95 often remains under 1\,s per chunk; the
$<3$\,s/chunk product target is evaluated with GPU SepFormer and
\texttt{Whisper-small} (beam size\,2) as in the reference configuration.
Detailed per-component latency breakdown is in
\texttt{outputs/metrics/latency\_benchmark.csv}.

\subsection{Error Analysis}

Error analysis (\texttt{notebooks/07\_error\_analysis.ipynb}) suggests
the following patterns on real multi-channel data (qualitative; exact
rates depend on session and road noise):
\begin{itemize}
    \item \textbf{Noise robustness}: WER tends to increase for
    sessions recorded at highway speed (lower cabin SNR).
    \item \textbf{Utterance length}: Short utterances ($<$2\,s) show higher
    deletion rates due to Whisper's minimum-context requirement.
    \item \textbf{Speaker overlap}: Sessions with simultaneous speech (true
    overlap, not just turn-taking) present the hardest separation challenge.
\end{itemize}

% ============================================================
% Section 5: Conclusion & Future Work
% ============================================================

\section{Conclusion}
\label{sec:conclusion}

\subsection{Summary}

We presented InCar-MSAR, the first end-to-end, open-source pipeline for
multi-speaker automatic speech recognition in automotive cabin environments.
Our three main contributions are: (1) the first published benchmark on
AISHELL-5, establishing WER and cpWER baselines; (2) empirical demonstration
that a separation-first pipeline (SepFormer + Whisper) substantially
outperforms single-channel direct transcription; and (3) a reproducible,
containerized system with a live Streamlit demo ready for investor and OEM
presentations. Experimental results on 100 sessions from AISHELL-5 eval1
confirm that each pipeline component contributes positively to overall
accuracy (see ablation in Table~\ref{tab:ablation}).

\subsection{Limitations}

Despite the encouraging results, the current system has notable limitations:

\begin{itemize}
    \item \textbf{Language scope}: The system is tuned for Mandarin Chinese
    (AISHELL-5). Extension to multilingual in-car ASR (e.g., Vietnamese,
    Japanese) would require retraining or adapting the intent keyword mapping.
    \item \textbf{Separation model mismatch}: SepFormer was pretrained on
    WSJ0-2Mix (clean, anechoic mixtures) and is applied zero-shot to
    reverberant car-cabin audio. Fine-tuning on AISHELL-5 training data
    (near-field references available) is expected to yield significant gains.
    \item \textbf{Speaker count flexibility}: The pipeline assumes $N{=}2$
    speakers by default; sessions with 3–4 simultaneous speakers degrade
    separation quality and require explicit configuration.
\end{itemize}

\subsection{Future Work}

Several promising directions emerge from this study:

\begin{enumerate}
    \item \textbf{Domain-adaptive fine-tuning}: Fine-tune SepFormer and
    Whisper jointly on AISHELL-5 near-field references (available in the
    54\,GB training split) to reduce the train-test domain gap.
    \item \textbf{On-device deployment}: Quantize the pipeline (INT8/INT4)
    for embedded deployment on automotive-grade SoCs (e.g., Qualcomm SA8295P).
    \item \textbf{LLM-based intent understanding}: Replace the keyword engine
    with a small instruction-tuned LLM (e.g., Qwen-1.5-1.8B) for richer
    intent extraction and multi-turn dialogue.
    \item \textbf{Real-time hardware integration}: Integrate with
    vehicle CAN bus simulators to enable actuator-in-the-loop evaluation.
    \item \textbf{Multilingual extension}: Leverage Whisper's multilingual
    backbone to extend the demo to Vietnamese and English cabin scenarios.
\end{enumerate}

\subsection{Broader Impact}

The open-source release of InCar-MSAR benefits the research community by
providing a reproducible baseline on AISHELL-5 and lowering the barrier to
in-car multi-speaker ASR research. Commercial applications in OEM
infotainment and ride-sharing services could improve accessibility for
hearing-impaired passengers and enhance safety by reducing driver distraction.
All third-party assets are used in accordance with their respective licenses
(AISHELL-5: CC BY-SA 4.0; Whisper: MIT; SpeechBrain: Apache 2.0).


% ─── Acknowledgments ───────────────────────────────────────
\section*{Acknowledgments}
We thank the AISHELL team for releasing AISHELL-5 under CC BY-SA 4.0.
This work uses Whisper (MIT License), SpeechBrain (Apache 2.0), and Asteroid (MIT License).

% ─── References ────────────────────────────────────────────
\bibliographystyle{plainnat}
\bibliography{references}

\end{document}
